\section{Data and Measurement}
\label{sec:data}

This section describes the price data underlying the analysis, the construction of the
realized-volatility outcome, a calendar-grid correction that materially affects inference, the
choice of donor (control) commodities, and the principled exclusions we impose before estimation.
We also state the data limitations honestly: there are no post-suspension futures for the banned
commodities, and the district-level food-staple donor pool is still being assembled at the time of
writing.

\subsection{Source: CEDA / MoA Agmarknet mandi spot prices}

Our primary data are daily wholesale (\emph{mandi}) spot prices for agricultural commodities,
sourced through the Centre for Economic Data and Analysis (CEDA) Data Portal API, which re-serves
the official Agmarknet price archive of India's Ministry of Agriculture and Farmers' Welfare (MoA).
Throughout we cite both the proximate provider (CEDA) and the originating authority (MoA
Agmarknet); the API delivers cleaned, documented, programmatically retrievable series that we
validate against the public Agmarknet portal. The reported price is the daily modal (most-frequent)
transaction price by market, the standard Agmarknet spot concept used in the related literature on
the suspension (e.g., Aggarwal, Chatterjee \& Sehgal 2023). The CEDA API enforces a hard rate limit
(40 requests/hour), so the disaggregated district pull is run as a paced, resumable job rather than
a single bulk download.

The sample spans 2017-01-02 through 2025-10-30 at daily frequency. We retrieve data at two levels of
aggregation: (i) a national series per commodity, used for headline charts only, and (ii) a
district-level panel, which is the basis for all estimation. The full district pull comprises on the
order of 5.74 million market-day observations across roughly 1{,}994 commodity--district units.

We deliberately use spot (mandi) prices rather than futures for the volatility analysis. The
suspension by construction halted futures trading in the banned commodities, so a post-period
futures series for the treated group does not exist (Section~\ref{subsec:limits}); spot prices,
however, are quoted continuously across the entire window for both treated and control commodities,
making them the only series on which a pre/post, treated/control comparison is even defined. Crude
palm oil (CPO) is the exception among the banned set: it has no representative mandi spot series and
is therefore routed out of the spot-volatility track entirely (toward the domestic--international
spread analysis discussed under~H4), not included among the treated units here.

\subsection{The district panel}

From the daily district-level prices we construct a monthly panel of realized volatility indexed by
commodity $\times$ district $\times$ month. The assembled panel contains $130{,}143$
commodity--district--month observations over 2017-02 through 2025-10, spanning $1{,}994$
commodity--district units; after listwise deletion of rows with missing leads/lags and the
within-window exclusions, the \emph{headline DiD is estimated on} $104{,}491$ observations across
$1{,}570$ commodity--district units (the figure reported in Table~\ref{tab:robustness}).
Working at the district level is a deliberate
identification choice. Treatment is assigned at the \emph{commodity} level by the regulator, so the
panel has only a handful of distinct treated clusters (five, after the exclusions below); the
district disaggregation supplies genuine cross-sectional variation and statistical power for
panel and synthetic-control estimators, while we continue to \emph{cluster all inference at the
commodity level} so as never to over-credit the district count as if it were independent
information. National-mean series are reserved for descriptive charts only: they are not robust (a
single $\sim$\,\rupee14.3M guar outlier distorts the national mean, which is why district series are
built as medians) and reduce to a coarse few-versus-few comparison.

\subsection{The realized-volatility measure}

Our outcome is monthly realized volatility computed from daily spot returns. For each
commodity--district series we form daily log returns, compute the 30-day rolling standard deviation
of those returns, and annualize by the conventional $\sqrt{252}$ trading-day factor; we then sample
this series monthly. The estimation outcome is the natural log of this 30-day annualized realized
volatility, $\ln(\mathrm{rv30})$, which stabilizes the variance and gives the regression
coefficients a percentage-change interpretation. Pre-registered robustness twins use a 60-day window
in place of the 30-day window, and a Parkinson high--low estimator where OHLC spot data are
available.

\subsection{The trading-day (calendar-grid) correction and why it matters}

A subtle but consequential data-handling issue concerns the calendar on which returns are computed.
The raw mandi series, as delivered, behaved as a 7-day calendar grid: weekend dates carried prices
forward, so the series contained Saturday and Sunday observations even though no trading occurs on
those days. Computing daily returns on that 7-day grid and then annualizing with the
\emph{trading-day} factor $\sqrt{252}$ is a units mismatch---the return interval and the
annualization constant assume different calendars. The fix is to restrict every series to weekdays
(Monday--Friday) \emph{before} forming returns and realized volatility, so that the return spacing
and the $\sqrt{252}$ annualization are consistent.

This correction is not cosmetic; it changes the qualitative verdict of the falsification battery.
On the uncorrected (calendar-grid) data the naive difference-in-differences design \emph{failed} its
integrity checks: a fake-ban placebo dated to December~2019 spuriously ``found'' an effect
($p \approx 0.046$), and the joint pre-trend lead test \emph{rejected} parallel trends
($p \approx 0.033$). On those grounds the design looked dead. After the weekday filter, dropping
MSP-censored paddy, and expanding the donor pool toward non-banned food commodities, the same
diagnostics are much cleaner: the December-2019 placebo is small and insignificant
($-6.5\%$; analytic $p \approx 0.344$, wild-cluster bootstrap $p \approx 0.433$), and the joint
pre-trend test no longer rejects parallel trends ($p \approx 0.445$). The current headline estimate
is correspondingly more conservative than the earlier clean-core run: the aggregate DiD is about
$-9.8\%$ with few-cluster inference above conventional significance thresholds. The lesson we draw,
and report, is that a mechanical calendar artifact had been driving the apparent failure of the
design, while donor choice materially affects the aggregate magnitude; the corrected food-donor
series are the basis for everything that follows.

\subsection{Donor (control) commodities: Option~B}

Synthetic-control and difference-in-differences estimators require untreated commodities to build the
counterfactual. Following the precedent set by the related suspension literature (Aggarwal,
Chatterjee \& Sehgal 2023; Gaurav \& Pandey 2024), donors must satisfy three criteria: they should
belong to the same or an adjacent commodity group where possible, they must have remained traded and
liquid through the suspension, and they must carry no own treatment (no commodity-specific policy
shock) within the estimation window.

We adopt what we label the ``Option~B'' donor pool. The clean, screened core, already in hand, is
\{castor, guarseed413, cotton, jeera, turmeric\}. To this core we are adding non-banned food and
oilseed staples---barley, maize, jowar, bajra, ragi, groundnut, sesamum, and sunflower---which
restore the food-inflation / MSP channel that the treated cereals and oilseeds (chana, wheat,
mustard) load on, and which the industrial/spice core does not span. The clean five are sufficient to
produce point estimates today; the food staples are being incorporated into the district panel to
sharpen inference (see below). Policy-ledger exclusion windows are applied before estimation---for
example, cotton's August-2022 unit break is confined to the futures track and does not contaminate
the spot donor series.

\subsection{Principled exclusions}

Three data-quality exclusions are imposed before estimation, each pre-registered.

\paragraph{Paddy dropped (MSP censoring).}
Paddy is dropped from the primary analysis because its spot price is effectively
\emph{censored} by minimum-support-price (MSP) procurement. Government procurement (FCI / state
agencies) pins the paddy price at the support level for extended stretches: $40.3\%$ of paddy's daily
spot returns are \emph{exactly} flat, against $1$--$6\%$ for clean, market-traded commodities.
Realized volatility computed on an administratively pinned price is a mechanical artifact of the
support regime, not a measure of market volatility, so paddy cannot inform a volatility-effect
estimate and is excluded. Wheat is also MSP-exposed ($19.7\%$ flat returns) but is one of the three
core commodities of interest; we therefore \emph{retain} wheat with an explicit MSP flag and read its
result as a robustness check rather than as a clean treated unit.

\paragraph{Guar identifier fix.}
The guar control required a vendor-identifier correction. The guar series under identifier~75 is
gum-contaminated and correlates only $0.04$ with the corresponding futures underlying; the clean guar
underlying (\texttt{guarseed413}, identifier~413) correlates $0.99$ and is the series we use as the
control. Identifier~75 is dropped.

\subsection{Data limitations}
\label{subsec:limits}

We are explicit about what the data cannot do. First, \emph{there are no post-suspension futures for
the banned commodities}: the suspension itself halted that trading, so any analysis requiring a
post-period banned-commodity futures series---most importantly a conventional futures--spot basis---is
infeasible as originally posed. Of the seven banned commodities, only CPO (traded on MCX) yields any
post-suspension banned-futures data; the six NCDEX-traded banned futures are unavailable. This is a
binding data limitation, not a modeling choice, and it is the reason the basis arm is reframed toward
a domestic-versus-international spread rather than a domestic futures--spot basis. Second, the
district-level pull of the food-staple donors is still in progress; the current estimates rest on the
clean industrial/spice core, and incorporating the food donors at the district level is the explicit
route to widening the donor pool and sharpening the few-treated-cluster inference. Finally, inference
remains coarse---roughly five treated commodity clusters---and the wheat result must be read against
its MSP exposure. We therefore frame the headline as a point estimate whose credibility rests on its
stability to leaving out any single treated commodity, not on the across-estimator agreement (which,
sharing one panel and five treated/five donor commodities, is largely mechanical), and as a
quasi-causal / descriptive association rather than a fully identified causal effect.
